% ============================================================
% PHỤ LỤC: CODE SNIPPETS
% ============================================================
\chapter{Phụ lục: Code nguồn chính}
\label{ch:appendix}

\section{Kiến trúc YOLO26n}

\begin{lstlisting}[caption={build\_yolo26n --- Xây dựng model YOLO26 scale nano}]
def build_yolo26n(nc=80):
    d, w, mc = 0.50, 0.25, 1024
    def ch(c): return max(round(min(c, mc) * w), 1) & ~7
    def rep(n): return max(round(n * d), 1)

    backbone = nn.ModuleList([
        Conv(3,      ch(64),  3, 2),      # 0 P1/2
        Conv(ch(64), ch(128), 3, 2),      # 1 P2/4
        C3k2(ch(128),ch(256),rep(2),False,0.25),
        Conv(ch(256),ch(256), 3, 2),      # 3 P3/8
        C3k2(ch(256),ch(512),rep(2),False,0.25),
        Conv(ch(512),ch(512), 3, 2),      # 5 P4/16
        C3k2(ch(512),ch(512),rep(2),True),
        Conv(ch(512),ch(1024),3, 2),      # 7 P5/32
        C3k2(ch(1024),ch(1024),rep(2),True),
        SPPF(ch(1024),ch(1024),5,3),
        C2PSA(ch(1024),ch(1024),rep(2)),
    ])
    detect = Detect(nc=nc, reg_max=1, end2end=True,
                    ch=(ch(256), ch(512), ch(1024)))
    return backbone, head, detect
\end{lstlisting}

\section{make\_anchors}

\begin{lstlisting}[caption={make\_anchors --- Tạo anchor points cho detection}]
def make_anchors(feats, strides, grid_cell_offset=0.5):
    anchor_points, stride_tensor = [], []
    for feat, stride in zip(feats, strides):
        _, _, h, w = feat.shape
        sx = torch.arange(w, device=feat.device) + grid_cell_offset
        sy = torch.arange(h, device=feat.device) + grid_cell_offset
        sy, sx = torch.meshgrid(sy, sx, indexing='ij')
        anchor_points.append(torch.stack([sx, sy], -1).view(-1, 2))
        stride_tensor.append(torch.full((h*w, 1), stride, device=feat.device))
    return torch.cat(anchor_points), torch.cat(stride_tensor)
\end{lstlisting}

\section{dist2bbox và bbox2dist}

\begin{lstlisting}[caption={dist2bbox --- Chuyển LTRB distances thành bounding box}]
def dist2bbox(distance, anchor_points, xywh=True, dim=-1):
    lt, rb = distance.chunk(2, dim)
    x1y1 = anchor_points - lt
    x2y2 = anchor_points + rb
    if xywh:
        c_xy = (x1y1 + x2y2) / 2
        wh   = x2y2 - x1y1
        return torch.cat([c_xy, wh], dim)
    return torch.cat([x1y1, x2y2], dim)

def bbox2dist(anchor_points, bbox, reg_max=None):
    x1y1, x2y2 = bbox.chunk(2, -1)
    dist = torch.cat((anchor_points - x1y1, x2y2 - anchor_points), -1)
    if reg_max is not None:
        dist = dist.clamp_(0, reg_max - 0.01)
    return dist
\end{lstlisting}

\section{Newton--Schulz Iteration}

\begin{lstlisting}[caption={Newton--Schulz orthogonalization trong Muon optimizer}]
def newton_schulz_orthogonalize(G, steps=5):
    # Chuan hoa Frobenius
    X = G / G.norm()  # ||X_0||_F = 1

    # He so cubic convergence
    a, b, c = 3.4445, -4.7750, 2.0315

    for _ in range(steps):
        A = X @ X.T
        X = a * X + b * (A @ X) + c * (A @ (A @ X))

    return X  # X ~ U V^T
\end{lstlisting}

\section{E2ELoss decay}

\begin{lstlisting}[caption={E2ELoss --- Progressive decay O2M/O2O weights}]
class E2ELoss:
    def __init__(self, model):
        self.o2m = 0.8
        self.o2o = 0.2
        self.final_o2m = 0.3

    def decay(self, x):
        return (max(1 - x / max(epochs - 1, 1), 0)
                * (self.o2m_copy - self.final_o2m)
                + self.final_o2m)

    def update(self):
        self.updates += 1
        self.o2m = self.decay(self.updates)
        self.o2o = max(1.0 - self.o2m, 0)
        # ProgLoss: box/cls weight update
        progress = min(self.updates / max(self.prog_epochs - 1, 1), 1.0)
        self.one2many.box_weight = 7.5 + (3.75 - 7.5) * progress
        self.one2many.cls_weight = 0.5 + (1.5 - 0.5) * progress
\end{lstlisting}
